\documentclass{article}

% Pass 'final' to the neurips_2023 package to remove line numbers if needed
\usepackage[final]{neurips_2023}

\usepackage{graphicx}

\usepackage[utf8]{inputenc} % allow utf-8 input
\usepackage[T1]{fontenc}    % use 8-bit T1 fonts
\usepackage{hyperref}       % hyperlinks
\usepackage{url}            % simple URL typesetting
\usepackage{booktabs}       % professional-quality tables
\usepackage{amsfonts}       % blackboard math symbols
\usepackage{nicefrac}       % compact symbols for 1/2, etc.
\usepackage{microtype}      % microtypography
\usepackage{xcolor}         % colors

\usepackage{hyperref}
\usepackage{cleveref}

\title{Summarization: Learning a Generalizable Trajectory Sampling Distribution for Model Predictive Control}

% Paper BibTex
% @article{power2024learning,
%   title={Learning a generalizable trajectory sampling distribution for model predictive control},
%   author={Power, Thomas and Berenson, Dmitry},
%   journal={IEEE Transactions on Robotics},
%   volume={40},
%   pages={2111--2127},
%   year={2024},
%   publisher={IEEE}
% }

\author{
  Xiaonan (Nice) Wang (wangxiaonannice@gmail.com) \\
  % Paper fr last-names (year)
  Summarization Generated based on Paper from Power and Berenson (2024) \\
  Topic: MPC, Generalizable Trajectory Sampling, Collision-Free Navigation, Prior \& Posterior Flow \\
}

\begin{document}

\maketitle

\begin{abstract}
This document summarizes the core contributions and methodology of the paper "Learning a Generalizable Trajectory Sampling Distribution for Model Predictive Control, \cite{power2024learning}", focusing on its' main ideas and the core blocks.
\end{abstract}

\section{Core Questions and Answers}

\paragraph{(1) What is the problem?} 
\begin{itemize}
  \item \textbf{Collision-Free Navigation} in Cluttered Environments
  \item Learning \textbf{Generalizable} Trajectories in \textbf{Out-of-Distribution (OOD)} Environments
\end{itemize}

\paragraph{(2) Why need to solve this problem?} 
Robotics applications like autonomous driving or manipulation require rapid and reliable trajectory generation. In complex environments, random sampling of control sequences is unlikely to generate low-cost paths, causing robots to get stuck in local optimal. Bridging the gap between sim-to-real models and diverse real-world (eg. OOD environments) scenarios is critical for robust autonomous operation.

\paragraph{(3) How is it different from previous methods?} 
\begin{itemize}
  \item Unlike previous work that utilizes simple Gaussian priors or uninformed sampling distributions, this approach employs a \textbf{Conditional Normalizing Flow (CNF)} to learn a goal-directed and environment-aware control sequence \textbf{Posterior}.
  \item It is also a sampling-based MPC framework to learn this \textbf{Posterior} from a dataset.
  \item Furthermore, it introduces a novel \textbf{Environment Projection} method to adapt to OOD scenarios by converting them to in-distribution representations during online navigation.
\end{itemize}

\paragraph{(4) Why is it better than previous methods? (Advantages)} 
\begin{itemize}
  \item \textbf{General Architecture}: The learned distribution can be integrated into different MPC frameworks (FlowMPPI, FlowiCEM).
  \item \textbf{Adaptability}: The projection mechanism allow effective generalizable trajectory sampling to OOD real-world environments.
  \item \textbf{Expressiveness}: It captures multimodal (conditional) trajectories which standard Gaussian-based samplers cannot.
\end{itemize}

\begin{figure}[tb]
  \centering
  \includegraphics[width=0.95\textwidth]{fig/flow_mpc_diagram.png}
  \caption{FlowMPC Block-Diagram (Mainly Inference): Use FlowMPPI as Example}
  \label{fig:flow_mpc_architecture}
\end{figure}

\paragraph{(5) What is the approach itself?} 
\begin{itemize}
  \item FlowMPC (FlowMPPI \& FlowiCEM)
  \begin{itemize}
    \item Normalizing Flow:
      \begin{itemize}
        \item Prior Flow
        \item Posterior Flow for Sampling
      \end{itemize}
    \item MPC (MPPI \& iCEM)
    \item Projection: Generalizing to Out-of-Distribution (OOD) Environments
  \end{itemize}
\end{itemize}

\paragraph{(6) What are the applications of it?} 
The approach can be demonstrated on several robotic platforms for collision-free navigation and manipulation tasks:
\begin{itemize}
  \item 2-D double-integrator
  \item 12-DoF quadrotor
  \item 7-DoF kinematic manipulator
  \item Typically, most robotic manipulators and navigators (eg. for vehicles)
\end{itemize}

% \newpage

\section{The Structure}

\paragraph{Summarized Block-Diagram} 
\begin{itemize}
  \item Mainly Inference
  \item Use FlowMPPI as Example
\end{itemize}
The summarized block-diagram see \cref{fig:flow_mpc_architecture}.

\paragraph{Original in Paper}
The original block-diagram see \cref{fig:original}.

\begin{figure}[tb]
  \centering
  \includegraphics[width=0.95\textwidth]{fig/orig.png}
  \caption{Original Block-Diagram in Paper}
  \label{fig:original}
\end{figure}

\bibliographystyle{abbrvnat}
\bibliography{main}
\end{document}